\documentclass[11pt]{article}
\usepackage[margin=1in]{geometry}
\usepackage{amsmath,amssymb,amsthm,mathtools}
\usepackage{booktabs}

\newcommand{\Gt}{\mathbb{G}_T}
\newcommand{\Gone}{\mathbb{G}_1}
\newcommand{\Gtwo}{\mathbb{G}_2}
\newcommand{\Fr}{\mathbb{F}_r}
\newcommand{\Fp}{\mathbb{F}_p}
\newcommand{\Fpk}{\mathbb{F}_{p^{12}}}
\newcommand{\MLpre}{\mathrm{ML_{pre}}}
\newcommand{\FE}{\mathrm{FE}}
\newcommand{\midd}{\mathrm{mid}}

\theoremstyle{plain}
\newtheorem{theorem}{Theorem}
\newtheorem{lemma}[theorem]{Lemma}
\newtheorem{proposition}[theorem]{Proposition}

\theoremstyle{remark}
\newtheorem*{remark}{Remark}

\title{Middle Product via Pairings on BLS12-381}
\author{}
\date{}

\begin{document}
\maketitle

%======================================================================
\part*{I.\quad Polynomial algorithms}
%======================================================================

\medskip\noindent\textit{The algorithms in this part are due to
Hanrot, Quercia, and Zimmermann~\cite{HQZ04}.
We give self-contained proofs for the reader's convenience.}

\section{Problem}

Let $R$ be a commutative ring.
Given polynomials
$A(x) = \sum_{i=0}^{N-1} a_i x^i$
and $B(x) = \sum_{j=0}^{2N-2} b_j x^j$
over~$R$,
their product $C = A \cdot B$ has degree $\le 3N{-}3$,
with coefficients
$c_k = \sum_{i+j=k} a_i b_j$.

\medskip\noindent\textbf{Goal.}
Compute the $N$ \emph{middle coefficients}
$c_{N-1}, c_N, \ldots, c_{2N-2}$.

\medskip
We write $\midd(A, B)$ for the vector of middle coefficients:
\begin{equation}\label{eq:mid}
\midd(A,B)_m = c_{N-1+m} = \sum_{i=0}^{N-1} a_i\, b_{N-1+m-i},
\qquad m = 0, \ldots, N{-}1.
\end{equation}

\noindent
In many applications $B$ is fixed and only~$A$ varies.
This allows precomputing $B$-dependent data
to reduce the per-query cost;
we discuss this optimization separately
in each algorithm below.

\section{Karatsuba middle product}

\subsection{Base case ($N = 2$)}

Let $A = (a_0, a_1)$ and $B = (b_0, b_1, b_2)$.
The middle coefficients are
$c_1 = a_0 b_1 + a_1 b_0$
and $c_2 = a_0 b_2 + a_1 b_1$.

\medskip\noindent
The naive computation uses four multiplications.
Using only three:
\begin{align*}
\alpha &= a_1 \cdot (b_0 + b_1), \\
\beta &= (a_1 - a_0) \cdot b_1, \\
\gamma &= a_0 \cdot (b_1 + b_2).
\end{align*}
Then:
\begin{align*}
\alpha - \beta &= a_1 b_0 + a_1 b_1 - a_1 b_1 + a_0 b_1
= a_0 b_1 + a_1 b_0 = c_1, \\
\gamma + \beta &= a_0 b_1 + a_0 b_2 + a_1 b_1 - a_0 b_1
= a_1 b_1 + a_0 b_2 = c_2.
\end{align*}
Each of the three multiplications pairs
a linear form in~$A$ with a linear form in~$B$.
If $B$ is fixed, the $B$-dependent quantities
$(b_0 + b_1)$, $b_1$, $(b_1 + b_2)$
can be precomputed.

\subsection{Splitting identity and algorithm ($N = 2^k$)}

Assume $N = 2^k \ge 2$ and split into equal halves:
\[
A_{\mathrm{lo}} = (a_0, \ldots, a_{N/2-1}), \qquad
A_{\mathrm{hi}} = (a_{N/2}, \ldots, a_{N-1}).
\]
Define three overlapping windows of~$B$,
each of length $N{-}1$:
\begin{align*}
B^{(0)} &= (b_0, \ldots, b_{N-2}), \\
B^{(1)} &= (b_{N/2}, \ldots, b_{N/2+N-2}), \\
B^{(2)} &= (b_N, \ldots, b_{2N-2}).
\end{align*}

\begin{proposition}\label{prop:split}
The middle coefficients decompose as:
\begin{align*}
\midd(A, B)_m &=
\midd(A_{\mathrm{lo}}, B^{(1)})_m
+ \midd(A_{\mathrm{hi}}, B^{(0)})_m,
\quad m = 0, \ldots, \tfrac{N}{2}{-}1, \\[3pt]
\midd(A, B)_{N/2+m} &=
\midd(A_{\mathrm{lo}}, B^{(2)})_m
+ \midd(A_{\mathrm{hi}}, B^{(1)})_m,
\quad m = 0, \ldots, \tfrac{N}{2}{-}1.
\end{align*}
\end{proposition}

\begin{proof}
From~\eqref{eq:mid},
$c_{N-1+m} = \sum_{i=0}^{N-1} a_i\, b_{N-1+m-i}$.
Splitting at $i = N/2$:
\[
c_{N-1+m}
= \underbrace{\sum_{i=0}^{N/2-1} a_i\, b_{N-1+m-i}}_{
  \text{$b$-indices in } [N/2{+}m,\; N{-}1{+}m]}
\;+\;
\underbrace{\sum_{j=0}^{N/2-1} a_{N/2+j}\, b_{N/2-1+m-j}}_{
  \text{$b$-indices in } [m,\; N/2{-}1{+}m]}.
\]
For $0 \le m < N/2$:
the first sum is $\midd(A_{\mathrm{lo}}, B^{(1)})_m$;
the second is $\midd(A_{\mathrm{hi}}, B^{(0)})_m$.
For $m = N/2 + m'$ with $0 \le m' < N/2$:
the first sum uses $B^{(2)}$;
the second uses $B^{(1)}$.
\end{proof}

\paragraph{Karatsuba's trick.}
The window $B^{(1)}$ appears in both output halves,
yielding four sub-middle-products.
We reduce to three.
Define
$S = B^{(0)} + B^{(1)}$,\;
$T = B^{(1)}$,\;
$U = B^{(1)} + B^{(2)}$,
and compute
\begin{alignat*}{2}
\alpha &= \midd(A_{\mathrm{hi}},\; S),\qquad
  &&\text{(size $N/2$)}\\
\beta  &= \midd(A_{\mathrm{hi}} - A_{\mathrm{lo}},\; T),\qquad
  &&\text{(size $N/2$)}\\
\gamma &= \midd(A_{\mathrm{lo}},\; U).\qquad
  &&\text{(size $N/2$)}
\end{alignat*}
By linearity of $\midd$ in both arguments:
\begin{align*}
\alpha - \beta
&= \midd(A_{\mathrm{hi}}, B^{(0)}{+}B^{(1)})
 - \midd(A_{\mathrm{hi}}{-}A_{\mathrm{lo}}, B^{(1)}) \\
&= \midd(A_{\mathrm{hi}}, B^{(0)})
 + \midd(A_{\mathrm{lo}}, B^{(1)})
= \text{lower half,} \\[3pt]
\gamma + \beta
&= \midd(A_{\mathrm{lo}}, B^{(1)}{+}B^{(2)})
 + \midd(A_{\mathrm{hi}}{-}A_{\mathrm{lo}}, B^{(1)}) \\
&= \midd(A_{\mathrm{hi}}, B^{(1)})
 + \midd(A_{\mathrm{lo}}, B^{(2)})
= \text{upper half.}
\end{align*}
Each recursive multiplication pairs
a linear form in~$A$ with a linear form in~$B$
($S$, $T$, $U$ depend only on~$B$).

\paragraph{Complexity.}
The recurrence $K(N) = 3\,K(N/2)$ with $K(1) = 1$
gives $K(N) = N^{\log_2 3} \approx N^{1.585}$ ring multiplications.
All other operations are additions.

\subsection{General algorithm (arbitrary $N$)}\label{sec:gen-alg}

The same splitting identity and Karatsuba trick
apply for arbitrary $N \ge 2$
with $N_0 = \lfloor N/2 \rfloor$,
$N_1 = \lceil N/2 \rceil = N - N_0$.
When $N$ is even, $N_0 = N_1$ and the algorithm
is identical to the power-of-two case.
The only change when $N$ is odd is the $\beta$~call:
$A_{\mathrm{hi}}$ has $N_1 = N_0 + 1$ entries while
$A_{\mathrm{lo}}$ has $N_0$, so we left-pad $A_{\mathrm{lo}}$
with a zero to form the difference
$Q = A_{\mathrm{hi}} - (0, A_{\mathrm{lo}})$.

\medskip\noindent
\textbf{Algorithm} $\mathrm{MP}(a[0\mathinner{.\,.} N{-}1],\;
  b[0\mathinner{.\,.} 2N{-}2])
  \;\to\; c[0\mathinner{.\,.} N{-}1]$.

\begin{enumerate}
\setcounter{enumi}{-1}
\item If $N = 1$: return $[a_0 \cdot b_0]$.
\item Set $N_0 \gets \lfloor N/2 \rfloor$,\;
      $N_1 \gets N - N_0$.

\item \emph{A-vectors.}
\begin{align*}
P[i] &\gets a_{N_0+i}, && i = 0,\ldots,N_1{-}1
  \quad\text{($= A_{\mathrm{hi}}$),} \\
R[i] &\gets a_i, && i = 0,\ldots,N_0{-}1
  \quad\text{($= A_{\mathrm{lo}}$).}
\end{align*}
If $N$ is even ($N_0 = N_1$):
\[
Q[i] \gets a_{N_0+i} - a_i, \qquad i = 0,\ldots, N_1{-}1.
\]
If $N$ is odd ($N_1 = N_0 + 1$):
\[
Q[0] \gets a_{N_0},\qquad
Q[i] \gets a_{N_0+i} - a_{i-1},\quad i = 1,\ldots, N_1{-}1.
\]

\item \emph{B-vectors.}
\begin{align*}
S[i] &\gets b_i + b_{N_1+i},
  && i = 0,\ldots, 2N_1{-}2
  \quad\text{(length $2N_1{-}1$),}\\
T[i] &\gets b_{N_1+i},
  && i = 0,\ldots, 2N_1{-}2
  \quad\text{(length $2N_1{-}1$),}\\
U[i] &\gets b_{N_1+i} + b_{2N_1+i},
  && i = 0,\ldots, 2N_0{-}2
  \quad\text{(length $2N_0{-}1$).}
\end{align*}

\item \emph{Three recursive calls.}
\begin{align*}
\alpha &\gets \mathrm{MP}(P,\; S),
  &&\text{size } N_1, \\
\beta  &\gets \mathrm{MP}(Q,\; T),
  &&\text{size } N_1, \\
\gamma &\gets \mathrm{MP}(R,\; U),
  &&\text{size } N_0.
\end{align*}

\item \emph{Recombine.}
\[
c_i \gets \alpha_i - \beta_i,\; i = 0,\ldots,N_1{-}1;
\qquad
c_{N_1+j} \gets \gamma_j + \beta_j,\; j = 0,\ldots,N_0{-}1.
\]
\end{enumerate}

\paragraph{Complexity.}
The multiplication count $K(N)$ satisfies
$K(1) = 1$ and
\[
K(N) = 2\,K(\lceil N/2 \rceil) + K(\lfloor N/2 \rfloor).
\]
This gives $K(N) = \Theta(N^{\log_2 3})$,
with $K(N) = N^{\log_2 3}$ exactly when $N$ is a power of two.
In the cost estimates of Part~II we write $K(N)$
for the exact count
and use $N^{1.585}$ as shorthand when $N = 2^k$.

\subsection{Precomputation for fixed $B$}

When $B$ is fixed across multiple queries,
the $B$-vectors $S$, $T$, $U$ at every recursion level
depend only on~$B$ and can be precomputed once.
The recursion tree has $K(N)$ leaves
(one ring element each).
If the data is laid out in depth-first traversal order,
the online phase streams through it sequentially.

When $B$ is not fixed, the $B$-vectors are computed
on-the-fly during each middle product evaluation.
The number of ring multiplications is unchanged;
only the additive overhead (forming $S$ and $U$)
is no longer amortized.

\section{FFT middle product}

\begin{lemma}\label{lem:cyclic}
Let $R_m$ denote the $m$-th coefficient of
the cyclic convolution
$A(x) \cdot B(x) \bmod (x^{2N} - 1)$.
Then the middle coefficients of $A \cdot B$
appear unmodified at positions $N{-}1$ through $2N{-}2$:
\[
R_m = c_m \quad \text{for all } m = N{-}1, \ldots, 2N{-}2.
\]
\end{lemma}

\begin{proof}
The product $A \cdot B$ has degree $\le 3N{-}3$,
so its coefficients $c_k$ are zero for $k \ge 3N{-}2$.
Reducing modulo $x^{2N} - 1$ wraps the coefficient at position $m + 2N$
onto position~$m$, giving
$R_m = c_m + c_{m+2N}$.
For the wraparound term $c_{m+2N}$ to vanish,
it suffices that $m + 2N > 3N{-}3$, i.e., $m > N{-}3$.
Since every middle position satisfies $m \ge N{-}1 > N{-}3$,
the wraparound vanishes and $R_m = c_m$.
\end{proof}

\paragraph{Algorithm.}
Let $N = 2^k$ and let $\omega \in R$
be a primitive $2N$-th root of unity.
\begin{enumerate}
\item DFT of $B$ (zero-padded to length~$2N$):
$\hat{B}_j = B(\omega^j)$ for $j = 0, \ldots, 2N{-}1$.
\item DFT of $A$ (zero-padded to length~$2N$):
$\hat{A}_j = A(\omega^j)$.
\item $2N$ pointwise multiplications:
$\hat{C}_j = \hat{A}_j \cdot \hat{B}_j$.
\item Inverse DFT: recover $R_0, \ldots, R_{2N-1}$.
\item Extract $R_{N-1}, \ldots, R_{2N-2}$.
\end{enumerate}

\noindent
When $B$ is fixed, step~1 is performed once
and reused for every new~$A$.

\paragraph{Cost.}
$2N$ multiplications in~$R$,
plus $O(N \log N)$ additions
and multiplications by roots of unity.
When $B$ is fixed, the online cost per query
is one forward DFT, $2N$ pointwise multiplications,
and one inverse DFT.

\section{Connection to bilinear complexity}

The Karatsuba and FFT middle product algorithms
can be understood through a more general principle.

\medskip
A \emph{bilinear algorithm} for a bilinear map
$\Psi(A, B) \colon R^p \times R^q \to R^r$
expresses it as
$\Psi(A,B) = \sum_{k=1}^{M} \alpha_k(A) \cdot \beta_k(B) \cdot w_k$,
where $\alpha_k$, $\beta_k$ are linear forms
and $w_k \in R^r$ are fixed.
The minimum such $M$ is the \emph{bilinear complexity}
of~$\Psi$.

Both the middle product and standard polynomial multiplication
are bilinear maps.
Fixing one argument (say~$A$) makes them linear,
and Tellegen's transposition principle~\cite{HQZ04}
shows that transposing a linear algorithm
for multiplication by~$A$
yields a linear algorithm for the middle product with~$A$ fixed,
preserving the number of multiplications.
By unfolding the dependence on~$A$,
this implies:

\begin{theorem}
The middle product
$\midd \colon R^N \times R^{2N-1} \to R^N$
and balanced polynomial multiplication
$\mathrm{Mult} \colon R^N \times R^N \to R^{2N-1}$
have the same bilinear complexity.
\end{theorem}

This gives an alternative proof that
$N^{\log_2 3}$ multiplications suffice
(via Karatsuba for balanced multiplication)
and that $2N$ multiplications suffice
(via FFT for balanced multiplication).
The explicit algorithms in Sections~2 and~3
can be seen as direct constructions achieving these bounds.

The proof and a detailed formal development
appear in the appendix of~\cite{HQZ04}.

\begin{thebibliography}{9}
\bibitem{HQZ04}
G.~Hanrot, M.~Quercia, and P.~Zimmermann,
\emph{The middle product algorithm~I: Speeding up the division
and square root of power series},
Appl.\ Algebra Eng.\ Commun.\ Comput.\ \textbf{14} (2004), 415--438.
\end{thebibliography}

%======================================================================
\part*{II.\quad Specialization to pairings on BLS12-381}
%======================================================================

\section{Setting}

We take $a_i \in \Gone$, $b_j \in \Gtwo$,
and ``multiplication'' $a_i \cdot b_j = e(a_i, b_j) \in \Gt$
via the BLS12-381 optimal ate pairing.
The output ring is $\Gt$ (written multiplicatively):
polynomial ``addition'' becomes $\Gt$~multiplication,
``subtraction'' becomes multiplication by an inverse.

\medskip\noindent
Throughout Part~II, $B$ is \textbf{fixed}
(i.e.\ the $\Gtwo$~points $b_0, \ldots, b_{2N-2}$ are known in advance).
For the naive and FFT methods,
we precompute $\Gtwo$~line data for the $O(N)$ fixed points;
the storage is modest.
For Karatsuba, we consider both
precomputed ($O(N^{1.585})$ storage) and
non-precomputed ($O(N)$ storage) variants.

\medskip\noindent
Two structural properties of~$\Gt$:
\begin{itemize}
\item \textbf{Cheap inversion.}
$\Gt$ lies in the cyclotomic subgroup of~$\Fpk^*$,
where $g^{-1} = g^{p^6}$:
a negation of one component, essentially free.

\item \textbf{The final exponentiation is a homomorphism.}
$f \mapsto f^{(p^{12}-1)/r}$ from~$\Fpk^*$ to~$\Gt$
commutes with multiplication and inversion.
\end{itemize}

\section{Cost model}

All costs are in \emph{pairing-equivalents}
($1 =$ one full pairing $\approx 0.75$\,ms with \texttt{blst}).

\begin{center}
\begin{tabular}{llr}
\toprule
\textbf{Operation} & \textbf{Symbol} & \textbf{Cost} \\
\midrule
Miller loop (precomputed $\Gtwo$ lines) & $\MLpre$ & 0.27 \\
Miller loop (no precomputation) & $\mathrm{ML}$ & 0.43 \\
Final exponentiation & $\FE$ & 0.57 \\
Full pairing with precomp ($\MLpre + \FE$) & & 0.84 \\
$\Gt$ exponentiation (255-bit, GLS-4) & & 0.50 \\
$\Gone$ scalar multiplication (255-bit, GLV) & & 0.10 \\
$\Fpk$ or $\Gt$ multiplication & & 0.04 \\
$\Gt$ inversion (cyclotomic conjugation) & & $\approx 0$ \\
\bottomrule
\end{tabular}
\end{center}

\noindent
The Miller loop / final exponentiation split
is from \texttt{blst} benchmarks.\footnote{%
\texttt{hackmd.io/@gnark/eccbench}}
The $\Gt$~exponentiation and $\Gone$~scalar multiplication costs
are estimates based on
GLS-4 and GLV algorithmic structure.

\section{Karatsuba with precomputation and deferred final exponentiation}

Here $B$ is fixed and we precompute
the entire tree of $B$-vectors from Section~2.3,
storing $\Gtwo$~line data at each of the $N^{1.585}$ leaves.

In the Karatsuba recursion, the recombination at each level
uses only $\Gt$~multiplication and inversion---no exponentiation.
Since the final exponentiation is a homomorphism,
it commutes past these operations.
We work with Miller loop outputs
(in~$\Fpk^*$) throughout the recursion,
applying the final exponentiation
only to the $N$ output values.

Cost per recursive pairing:
$\MLpre + 0.04 = 0.31$
(Miller loop plus $\Fpk$ recombination).
Plus $N$ final exponentiations at the end.

\[
\boxed{T_{\mathrm{kara}} = 0.31\,K(N) + 0.57\,N}
\]

\noindent
(For $N = 2^k$ this is $0.31\,N^{1.585} + 0.57\,N$.)

\paragraph{Storage.}
$K(N)$ leaves $\times$ ${\sim}\,20$\,KB line data each.
For $N = 2^k$ this is $N^{1.585} \times 20$\,KB.
Sequential access via depth-first layout.

\section{Karatsuba without precomputation (deferred final exponentiation)}

The $O(K(N))$ storage for the precomputed $B$-tree
can become prohibitive
(e.g.\ ${\sim}\,47$\,MB at $N = 128$,
${\sim}\,140$\,MB at $N = 256$).
If we forgo precomputation entirely,
each recursive pairing uses a full (unprecomputed) Miller loop
at cost $\mathrm{ML} = 0.43$ instead of $\MLpre = 0.27$.
The deferred final exponentiation still applies,
so the cost per recursive pairing becomes
$\mathrm{ML} + 0.04 = 0.47$.

\[
\boxed{T_{\mathrm{kara,nopre}} = 0.47\,K(N) + 0.57\,N}
\]

\noindent
(For $N = 2^k$ this is $0.47\,N^{1.585} + 0.57\,N$.)

\paragraph{Storage.}
$O(N)$ only (the input data).
No precomputation phase.

\paragraph{Trade-off.}
The ratio of leading coefficients is $0.47/0.31 \approx 1.52$,
so the no-precomp variant is about $52\%$ more expensive
per recursive pairing.
The advantage is zero setup cost
and dramatically less memory,
which may be decisive for large~$N$
or when $N$ is not known in advance.

\section{FFT with individual final exponentiations}

Since $B$ is fixed,
the $\Gtwo$-valued DFT outputs
$\hat{B}_j = B(\omega^j)$ are precomputed once.
For each~$\hat{B}_j$,
we also precompute the Miller loop line functions
(${\sim}\,20$\,KB per point using \texttt{blst\_precompute\_lines}),
reducing the online Miller loop cost
from $\mathrm{ML} = 0.43$ to $\MLpre = 0.27$ per pairing.

The inverse DFT butterflies
exponentiate $\Gt$~elements by twiddle factors
(${\sim}\,255$-bit elements of~$\Fr$),
using GLS-4 decomposition and cyclotomic squaring.
These require elements in the cyclotomic subgroup~$\Gt$,
i.e., fully exponentiated.
Therefore each of the $2N$ pairings
must include its own final exponentiation,
at a cost of $\MLpre + \FE = 0.84$ each.

Multi-pairing batching does not apply:
we need $2N$ individual $\Gt$ values, not their product.

With $L = \log_2(2N)$:
\[
\boxed{T_{\mathrm{fft}}
= \underbrace{1.68\,N}_{2N \text{ pairings at } 0.84}
\;+\; \underbrace{0.10\,NL}_{\Gone\text{ DFT}}
\;+\; \underbrace{0.54\,NL}_{\Gt\text{ IDFT}}
= 1.68\,N + 0.64\,NL}
\]

\paragraph{Storage.}
$2N$ precomputed $\Gtwo$ points with line data:
$O(N)$, about $40N$\,KB.

\section{Comparison}

\subsection{Cost formulas}

The naive method computes each output
as a product of $N$ pairings sharing one final exponentiation:
$T_{\mathrm{naive}} = 0.31\,N^2 + 0.57\,N$.
(The naive and FFT methods use $\Gtwo$~line precomputation
for the $O(N)$ fixed $B$-points,
which is modest in storage.)

\begin{center}
\begin{tabular}{rrrrrrr}
\toprule
$N$ & $L$ & Naive & Kara & Kara (no pre) & FFT & Best \\
\midrule
4 & 3 & 7 & 5 & 7 & 14 & Kara \\
8 & 4 & 24 & 13 & 17 & 34 & Kara \\
16 & 5 & 88 & 34 & 47 & 78 & Kara \\
32 & 6 & 336 & 94 & 132 & 177 & Kara \\
64 & 7 & 1{,}306 & 263 & 379 & 394 & Kara \\
128 & 8 & 5{,}152 & 751 & 1{,}101 & 870 & Kara \\
256 & 9 & 20{,}462 & 2{,}180 & 3{,}230 & 1{,}905 & FFT \\
512 & 10 & 81{,}557 & 6{,}395 & 9{,}543 & 4{,}137 & FFT \\
1{,}024 & 11 & 325{,}642 & 18{,}894 & 28{,}337 & 8{,}929 & FFT \\
\bottomrule
\end{tabular}
\end{center}

\subsection{Crossovers}

\paragraph{Naive vs.\ Karatsuba.}
Karatsuba (with or without precomputation) wins for all $N \ge 2$
(same leading coefficient, smaller exponent).

\paragraph{Karatsuba (precomp) vs.\ FFT.}
Around $N \approx 180$.
Two effects compound:

\emph{(1) Deferred FE helps Karatsuba.}
Karatsuba pays $\MLpre = 0.27$ per pairing,
amortizing $\FE = 0.57$ over $N$ outputs.

\emph{(2) Individual FE hurts FFT.}
The FFT pays $\MLpre + \FE = 0.84$ per pairing
because the $\Gt$~IDFT butterflies
require fully exponentiated inputs.

\emph{(3) Sensitivity to constants.}
Since $N^{0.585}$ grows slowly,
a factor-$\alpha$ change in the leading coefficient
shifts the crossover by~$\alpha^{1/0.585} \approx \alpha^{1.71}$.

\paragraph{Karatsuba (no precomp) vs.\ FFT.}
Around $N \approx 70$.
Without precomputation the Karatsuba leading coefficient
rises from $0.31$ to $0.47$,
so the FFT regime begins much sooner.
For applications where storage is constrained
or $N$ varies unpredictably,
this is the relevant crossover.

\subsection{Storage}

\begin{center}
\begin{tabular}{rrrrr}
\toprule
$N$ & Naive / FFT & Kara (no pre) & Kara ($\Gtwo$ only) & Kara (with lines) \\
\midrule
128 & ${\sim}\,5$\,MB & ${\sim}\,0$\,MB & ${\sim}\,0.5$\,MB & ${\sim}\,47$\,MB \\
256 & ${\sim}\,10$\,MB & ${\sim}\,0$\,MB & ${\sim}\,1.4$\,MB & ${\sim}\,140$\,MB \\
512 & ${\sim}\,20$\,MB & ${\sim}\,0$\,MB & ${\sim}\,4$\,MB & ${\sim}\,420$\,MB \\
\bottomrule
\end{tabular}
\end{center}

\noindent
``Kara (no pre)'' stores only the input data ($O(N)$).
Naive and FFT store precomputed $\Gtwo$~line data
for $O(N)$ points.

\section{Summary}

For $N$ in the range 100--1000:
\begin{itemize}
\item $N \lesssim 180$:
\textbf{Karatsuba with precomputation} is optimal,
if the $O(N^{1.585})$ storage is acceptable.
Deferred final exponentiation,
sequential memory access.

\item $N \lesssim 70$:
\textbf{Karatsuba without precomputation} is optimal
among methods with $O(N)$ storage.
Same algorithmic structure,
but each pairing uses a full Miller loop ($\mathrm{ML} = 0.43$).

\item $N \gtrsim 180$ (or $N \gtrsim 70$ without precomp):
\textbf{FFT} is optimal.
$O(N)$ storage, $O(N \log N)$ online cost.
Bottleneck is the $\Gt$~inverse DFT.

\item \textbf{Naive} is never optimal for $N \ge 2$.
\end{itemize}

\end{document}
